\begin{table}[!htbp]
\centering
\small
\renewcommand{\arraystretch}{1.3}
\begin{tabular}{>{\raggedright\arraybackslash}p{3.5cm}>{\raggedright\arraybackslash}p{2.6cm}>{\raggedright\arraybackslash}p{8.5cm}}
\toprule
\textbf{Mục tiêu an ninh} & \textbf{Vai trò của PKI} & \textbf{Cơ chế thực hiện / Phối hợp} \\
\midrule
\textbf{Xác thực (Authenticity)} & Hỗ trợ trực tiếp & Chứng thư số X.509, chữ ký số của CA trên Subject PublicKey. \\
\midrule
\textbf{Toàn vẹn (Integrity)} & Hỗ trợ trực tiếp & Hàm băm mật mã và chữ ký số trên cấu trúc chứng thư. \\
\midrule
\textbf{Chống chối bỏ (Non-repudiation)} & Hỗ trợ trực tiếp & Chữ ký số gắn liền với khóa bí mật duy nhất của thực thể. \\
\midrule
\textbf{Bí mật (Confidentiality)} & Kết hợp bên ngoài & Giao thức TLS bắt tay thỏa thuận khóa phiên (ECDHE) và mã hóa đối xứng (AES-GCM). \\
\midrule
\textbf{Sẵn sàng (Availability)} & Kết hợp bên ngoài & Hạ tầng mạng phân tán, bộ đệm chứng thực, OCSP Stapling, CDN. \\
\bottomrule
\end{tabular}
\caption[Phân định phạm vi hỗ trợ của PKI đối với mục tiêu an ninh thông tin]{Phân định phạm vi hỗ trợ của PKI đối với các mục tiêu an ninh thông tin.}
\label{tab:pki-security-goals}
\end{table}
